\documentclass[11pt,a4paper]{article}
\usepackage[top=2.3cm,bottom=2.3cm,left=2.2cm,right=2.2cm]{geometry}

\usepackage{amsmath,amssymb}
\usepackage{fontspec}
\setmainfont{Liberation Serif}
\setsansfont{Liberation Sans}
\setmonofont{DejaVu Sans Mono}[Scale=0.84]

\usepackage{xcolor}
\definecolor{brandblue}{RGB}{20, 55, 105}
\definecolor{accentblue}{RGB}{35, 95, 165}
\definecolor{codebg}{RGB}{248, 249, 251}
\definecolor{codeframe}{RGB}{215, 222, 230}
\definecolor{javakeyword}{RGB}{0, 45, 170}
\definecolor{javastring}{RGB}{5, 120, 20}
\definecolor{javacomment}{RGB}{120, 120, 120}
\definecolor{javanumber}{RGB}{20, 75, 220}

\definecolor{noteborder}{RGB}{30, 110, 65}
\definecolor{notebg}{RGB}{242, 249, 245}
\definecolor{warnborder}{RGB}{185, 75, 10}
\definecolor{warnbg}{RGB}{254, 248, 240}
\definecolor{tipborder}{RGB}{30, 95, 160}
\definecolor{tipbg}{RGB}{242, 247, 254}

\usepackage{fancyhdr}
\usepackage{lastpage}
\pagestyle{fancy}
\fancyhf{}
\lhead{\parbox[t]{0.58\textwidth}{\raggedright\small\sffamily\color{brandblue}\textbf{T10 --- Variabili, Scope, Call Stack e Memoria Heap}}}
\rhead{\small\sffamily\color{gray}Programmazione II -- UniVR (2025/2026)}
\cfoot{\small\sffamily Pagina \thepage\ di \pageref{LastPage}}
\renewcommand{\headrulewidth}{0.5pt}
\renewcommand{\headrule}{\hbox to\headwidth{\color{brandblue!70!black}\leaders\hrule height \headrulewidth\hfill}}

\usepackage{titlesec}
\titleformat{\section}{\Large\bfseries\sffamily\color{brandblue}}{\thesection}{0.8em}{}[\titlerule]
\titleformat{\subsection}{\large\bfseries\sffamily\color{accentblue}}{\thesubsection}{0.8em}{}
\titleformat{\subsubsection}{\normalsize\bfseries\sffamily\color{brandblue!85!black}}{\thesubsubsection}{0.8em}{}

\usepackage{needspace}
\usepackage{fancyvrb}
\DefineVerbatimEnvironment{asciibox}{Verbatim}{
  fontsize=\fontsize{7.5pt}{9.2pt}\selectfont,
  frame=single,
  rulecolor=\color{codeframe},
  fillcolor=\color{codebg},
  framesep=2.5mm
}
\usepackage{listings}
\lstset{
  backgroundcolor=\color{codebg},
  frame=single,
  rulecolor=\color{codeframe},
  basicstyle=\ttfamily\footnotesize,
  keywordstyle=\color{javakeyword}\bfseries,
  stringstyle=\color{javastring},
  commentstyle=\color{javacomment}\itshape,
  numberstyle=\tiny\color{gray},
  numbers=left,
  numbersep=7pt,
  stepnumber=1,
  breaklines=true,
  breakatwhitespace=true,
  showstringspaces=false,
  tabsize=4,
  xleftmargin=12pt,
  framexleftmargin=12pt
}

\newcommand{\passthrough}[1]{#1}
\providecommand{\tightlist}{\setlength{\itemsep}{0pt}\setlength{\parskip}{0pt}}

\usepackage{tcolorbox}
\tcbuselibrary{skins,breakable}
\newtcolorbox{studynote}[1][Nota]{
  enhanced,
  breakable,
  colback=notebg,
  colframe=noteborder,
  fonttitle=\bfseries\sffamily\small,
  coltitle=white,
  title={#1},
  arc=2mm,
  boxrule=0.8pt,
  left=9pt, right=9pt, top=7pt, bottom=7pt
}

\newtcolorbox{studywarning}[1][Attenzione]{
  enhanced,
  breakable,
  colback=warnbg,
  colframe=warnborder,
  fonttitle=\bfseries\sffamily\small,
  coltitle=white,
  title={#1},
  arc=2mm,
  boxrule=0.8pt,
  left=9pt, right=9pt, top=7pt, bottom=7pt
}

\newtcolorbox{studytip}[1][Suggerimento]{
  enhanced,
  breakable,
  colback=tipbg,
  colframe=tipborder,
  fonttitle=\bfseries\sffamily\small,
  coltitle=white,
  title={#1},
  arc=2mm,
  boxrule=0.8pt,
  left=9pt, right=9pt, top=7pt, bottom=7pt
}

\usepackage{booktabs}
\usepackage{longtable}
\usepackage{array}
\usepackage{calc}

\usepackage{hyperref}
\hypersetup{
  colorlinks=true,
  linkcolor=accentblue,
  urlcolor=accentblue,
  citecolor=brandblue,
  pdfborder={0 0 0}
}

\title{\Huge\bfseries\sffamily\color{brandblue} T10 --- Variabili, Scope, Call Stack e Memoria Heap \\[0.3em] \large\sffamily Stack Frame, Passaggio dei Parametri per Valore, Garbage Collection e Debutto di MavenDate}
\author{\textbf{Docente:} Prof. Michele Pasqua \quad---\quad \textbf{Appunti Revisione:} Wally}
\date{A.A. 2025/2026 -- Universit\`a degli Studi di Verona}

\begin{document}
\maketitle
\thispagestyle{fancy}
\vspace{-1.2em}
\noindent\textcolor{brandblue}{\rule{\textwidth}{0.8pt}}
\vspace{1.2em}

In questo blocco analizziamo in profondità il modello di memoria della
JVM (Stack vs Heap), il passaggio dei parametri per valore, gli oggetti
immutabili e i tipi wrapper
(\href{file:///home/wally/Documenti/GitHub/Programmazione-II/25-26/Teoria-e-Esercitazioni/Slides-20260902/T10\%20-\%20Java\%20Variables\%20and\%20Call\%20Stack.pdf}{T10
- Java Variables and Call Stack.pdf}). Sul fronte pratico analizziamo il
\textbf{grande spartiacque architetturale del corso}: il passaggio dal
progetto manuale a pacchetti
\href{file:///home/wally/Documenti/GitHub/Programmazione-II/25-26/Teoria-e-Esercitazioni/Codice-20260902/Lezione10/SimpleDate}{\passthrough{\lstinline!SimpleDate!}}
al progetto ufficiale
\textbf{\href{file:///home/wally/Documenti/GitHub/Programmazione-II/25-26/Teoria-e-Esercitazioni/Codice-20260902/Lezione10/MavenDate}{\passthrough{\lstinline!MavenDate!}}}
con Apache Maven, file
\href{file:///home/wally/Documenti/GitHub/Programmazione-II/25-26/Teoria-e-Esercitazioni/Codice-20260902/Lezione10/MavenDate/pom.xml}{\passthrough{\lstinline!pom.xml!}},
layout \passthrough{\lstinline!src/main/java!} e collisione di nomi tra
package diversi.

\begin{center}\rule{0.5\linewidth}{0.5pt}\end{center}

\subsubsection{1. Teoria: Concetti Teorici dalle Slide (Slide
T10)}\label{teoria-concetti-teorici-dalle-slide-slide-t10}

\paragraph{A. Il Modello di Memoria della JVM: Stack vs Heap (Slide
1--11)}\label{a.-il-modello-di-memoria-della-jvm-stack-vs-heap-slide-111}

La memoria a disposizione di un processo Java è partizionata in tre aree
principali: 1. \textbf{Call Stack (Stack dei Thread)}: - Memoria
estremamente veloce, gestita a politica LIFO (Push/Pop). - È suddivisa
in \textbf{Stack Frame} (o record di attivazione): ogni invocazione di
metodo crea un nuovo frame nello stack contenente: - La tabella delle
variabili locali. - I parametri formali passati al metodo. - Lo stack
degli operandi (per i calcoli intermedi della CPU/bytecode). - I
riferimenti alla costante del metodo e l'indirizzo di ritorno. - Quando
il metodo termina (con \passthrough{\lstinline!return!} o eccezione), il
frame viene rimosso istantaneamente dallo Stack e tutte le sue variabili
locali vengono distrutte. 2. \textbf{Heap (Memoria Dinamica degli
Oggetti)}: - Memoria globale condivisa da tutti i thread
dell'applicazione. - Ospita tutti gli \textbf{oggetti} istanziati con
\passthrough{\lstinline!new!}, inclusi gli \textbf{array}. - I dati
nello Heap non si deallocano all'uscita dal metodo: sopravvivono finché
esiste almeno un riferimento attivo nello Stack o in altri oggetti che
punta ad essi. Quando diventano irraggiungibili (\emph{unreachable}),
vengono eliminati dal Garbage Collector. 3. \textbf{Method Area /
Metaspace (Memoria di Classe)}: - Ospita il bytecode compilato delle
classi caricate dal ClassLoader, i metadati di tipo e le
\textbf{variabili statiche} (\passthrough{\lstinline!static!}). -
\textbf{Blocco di Inizializzazione Statica
(\passthrough{\lstinline!static \{ ... \}!})}: - Blocco speciale
eseguito \textbf{una sola volta} nel ciclo di vita dell'applicazione, al
momento in cui la classe viene caricata in memoria dalla JVM:
\passthrough{\lstinline!java     static int[] arr;     static \{         arr = new int[3];         arr[0] = 1; arr[1] = 2; arr[2] = 0;     \}!}
- Si usa per inizializzare strutture statiche complesse che richiedono
cicli o controlli prima dell'uso.

\begin{center}\rule{0.5\linewidth}{0.5pt}\end{center}

\paragraph{B. Immutabilità e Passaggio Parametri: Pass-by-Value (Slide
15--25)}\label{b.-immutabilituxe0-e-passaggio-parametri-pass-by-value-slide-1525}

\begin{itemize}
\tightlist
\item
  \textbf{Oggetti Immutabili (\emph{Immutable Objects})}:

  \begin{itemize}
  \tightlist
  \item
    Un oggetto il cui stato interno non può essere modificato dopo la
    costruzione (es. \passthrough{\lstinline!String!},
    \passthrough{\lstinline!Integer!}, \passthrough{\lstinline!Date!} da
    Lezione 09).
  \item
    \emph{Regole per creare una classe immutabile}:

    \begin{enumerate}
    \def\labelenumi{\arabic{enumi}.}
    \tightlist
    \item
      Rendere tutti i campi \passthrough{\lstinline!private final!}.
    \item
      Non fornire metodi mutatori (nessun setter).
    \item
      Dichiarare la classe \passthrough{\lstinline!final!} (per impedire
      che le sottoclassi violino l'immutabilità con override malevoli).
    \item
      Se l'oggetto contiene riferimenti a strutture mutabili (es. un
      array o un'altra data mutabile), eseguire copie difensive
      (\emph{defensive copy}) nel costruttore e nei getter.
    \end{enumerate}
  \item
    \emph{Vantaggi}: Immuni da side-effects da aliasing, intrinsecamente
    \emph{thread-safe} (possono essere condivisi tra thread concorrenti
    senza sincronizzazione o lock).
  \end{itemize}
\item
  \textbf{Come Java passa i parametri ai metodi: SEMPRE Strictly
  Pass-by-Value}:

  \begin{itemize}
  \tightlist
  \item
    In Java \textbf{non esiste il passaggio per riferimento del C++
    (\passthrough{\lstinline!int\& x!})}.
  \item
    \textbf{Per i tipi primitivi (\passthrough{\lstinline!int!},
    \passthrough{\lstinline!double!})}: il metodo riceve una copia del
    valore binario. Qualsiasi modifica locale sul parametro non ha alcun
    effetto sulla variabile del chiamante.
  \item
    \textbf{Per i tipi reference (oggetti/array)}: il metodo riceve
    \textbf{una copia per valore del riferimento} (dell'indirizzo
    puntato):

    \begin{itemize}
    \item
      \emph{Caso 1 (Mutazione dell'oggetto)}: se usiamo il riferimento
      ricevuto per chiamare un metodo mutatore
      (\passthrough{\lstinline!p.setAge(25)!}), l'oggetto puntato nello
      Heap viene effettivamente modificato per il chiamante!
    \item
      \emph{Caso 2 (Riassegnamento del puntatore --- Il tranello
      d'esame)}:

\begin{lstlisting}[language=Java]
void reset(Person p) {
    p = new Person("Bob"); // Assegna un nuovo oggetto al puntatore locale!
}
\end{lstlisting}

      All'esterno, la variabile del chiamante \textbf{non cambia
      affatto}! Ha semplicemente mutato la copia locale del puntatore
      nel proprio stack frame.
    \end{itemize}
  \end{itemize}
\end{itemize}

\begin{center}\rule{0.5\linewidth}{0.5pt}\end{center}

\paragraph{C. Classi Wrapper, Autoboxing e i loro Trabocchetti (Slide
26--35)}\label{c.-classi-wrapper-autoboxing-e-i-loro-trabocchetti-slide-2635}

\begin{itemize}
\tightlist
\item
  Per ciascuno degli 8 tipi primitivi, Java mette a disposizione una
  corrispondente \textbf{Wrapper Class} nel package
  \passthrough{\lstinline!java.lang!}:

  \begin{itemize}
  \tightlist
  \item
    \passthrough{\lstinline!byte!} \(\rightarrow\)
    \passthrough{\lstinline!Byte!}
  \item
    \passthrough{\lstinline!short!} \(\rightarrow\)
    \passthrough{\lstinline!Short!}
  \item
    \passthrough{\lstinline!int!} \(\rightarrow\)
    \passthrough{\lstinline!Integer!}
  \item
    \passthrough{\lstinline!long!} \(\rightarrow\)
    \passthrough{\lstinline!Long!}
  \item
    \passthrough{\lstinline!float!} \(\rightarrow\)
    \passthrough{\lstinline!Float!}
  \item
    \passthrough{\lstinline!double!} \(\rightarrow\)
    \passthrough{\lstinline!Double!}
  \item
    \passthrough{\lstinline!char!} \(\rightarrow\)
    \passthrough{\lstinline!Character!}
  \item
    \passthrough{\lstinline!boolean!} \(\rightarrow\)
    \passthrough{\lstinline!Boolean!}
  \end{itemize}
\item
  Tutte le classi wrapper sono \textbf{immutabili}.
\item
  \textbf{Autoboxing e Unboxing (Java 5+)}:

  \begin{itemize}
  \tightlist
  \item
    \emph{Autoboxing}: conversione automatica da primitivo a wrapper
    (\passthrough{\lstinline!Integer x = 5;!} \(\implies\) compilato
    come \passthrough{\lstinline!Integer.valueOf(5)!}).
  \item
    \emph{Unboxing}: estrazione automatica del valore primitivo dal
    wrapper (\passthrough{\lstinline!int y = x;!} \(\implies\) compilato
    come \passthrough{\lstinline!x.intValue()!}).
  \end{itemize}
\item
  Attenzione: \textbf{I Due Grandi Pericoli delle Wrapper Classes}:

  \begin{enumerate}
  \def\labelenumi{\arabic{enumi}.}
  \item
    \textbf{Il Trabocchetto dell'Integer Cache (\(-128 \dots +127\))}:

\begin{lstlisting}[language=Java]
Integer a = 100, b = 100;
System.out.println(a == b); // TRUE! (La JVM ricicla lo stesso oggetto cacheato)

Integer c = 200, d = 200;
System.out.println(c == d); // FALSE! (Fuori range [-128, 127], alloca 2 oggetti distinti nello Heap!)
\end{lstlisting}

    Conclusione: \textbf{Mai usare \passthrough{\lstinline!==!} per
    confrontare i wrapper}, usare sempre
    \passthrough{\lstinline!.equals()!}!
  \item
    \textbf{Crash da Unboxing su \passthrough{\lstinline!null!}}:

\begin{lstlisting}[language=Java]
Integer obj = null;
int val = obj; // RUNTIME ERROR: NullPointerException!
\end{lstlisting}

    La JVM tenta di invocare \passthrough{\lstinline!obj.intValue()!},
    fallendo rovinosamente se \passthrough{\lstinline!obj!} è nullo.
  \end{enumerate}
\end{itemize}

\begin{center}\rule{0.5\linewidth}{0.5pt}\end{center}

\subsubsection{\texorpdfstring{2. Codice: Evoluzione del Codice: Il
Passaggio a
\texttt{MavenDate}}{2. Codice: Evoluzione del Codice: Il Passaggio a MavenDate}}\label{codice-evoluzione-del-codice-il-passaggio-a-mavendate}

In
\href{file:///home/wally/Documenti/GitHub/Programmazione-II/25-26/Teoria-e-Esercitazioni/Codice-20260902/Lezione10}{\passthrough{\lstinline!Lezione10!}},
il docente archivia il vecchio progetto non strutturato e adotta lo
standard industriale \textbf{Apache Maven}.

\paragraph{1. Riorganizzazione dei
Package}\label{riorganizzazione-dei-package}

Il progetto viene suddiviso in due package distinti: -
\textbf{\passthrough{\lstinline!it.oop.core!}}: il dominio applicativo
(le classi di modello:
\href{file:///home/wally/Documenti/GitHub/Programmazione-II/25-26/Teoria-e-Esercitazioni/Codice-20260902/Lezione10/MavenDate/src/main/java/it/oop/core/Date.java}{\passthrough{\lstinline!Date.java!}},
\href{file:///home/wally/Documenti/GitHub/Programmazione-II/25-26/Teoria-e-Esercitazioni/Codice-20260902/Lezione10/MavenDate/src/main/java/it/oop/core/Language.java}{\passthrough{\lstinline!Language.java!}},
\href{file:///home/wally/Documenti/GitHub/Programmazione-II/25-26/Teoria-e-Esercitazioni/Codice-20260902/Lezione10/MavenDate/src/main/java/it/oop/core/BirthDay.java}{\passthrough{\lstinline!BirthDay.java!}}).
- \textbf{\passthrough{\lstinline!it.oop.ui!}}: l'interfaccia utente
contenente il \passthrough{\lstinline!main!}
(\href{file:///home/wally/Documenti/GitHub/Programmazione-II/25-26/Teoria-e-Esercitazioni/Codice-20260902/Lezione10/MavenDate/src/main/java/it/oop/ui/MainDate.java}{\passthrough{\lstinline!MainDate.java!}}).

\paragraph{2. Il Test Didattico sulla Collisione dei Nomi di
Package}\label{il-test-didattico-sulla-collisione-dei-nomi-di-package}

Nel package \passthrough{\lstinline!it.oop.ui!}, il docente aggiunge
appositamente una seconda classe:
\href{file:///home/wally/Documenti/GitHub/Programmazione-II/25-26/Teoria-e-Esercitazioni/Codice-20260902/Lezione10/MavenDate/src/main/java/it/oop/ui/Date.java}{\passthrough{\lstinline!it/oop/ui/Date.java!}}:

\begin{lstlisting}[language=Java]
package it.oop.ui;

class Date { // Visibilità package-private
    public int time;
    private int start = 0;

    public Date(int time) { this.time = time; }
}
\end{lstlisting}

E in
\href{file:///home/wally/Documenti/GitHub/Programmazione-II/25-26/Teoria-e-Esercitazioni/Codice-20260902/Lezione10/MavenDate/src/main/java/it/oop/ui/MainDate.java}{\passthrough{\lstinline!MainDate.java!}}
mostra come la JVM risolve l'ambiguità:

\begin{lstlisting}[language=Java]
package it.oop.ui;

import it.oop.core.Date; // 1. L'import esplicito ha la precedenza!

public class MainDate {
    public static void main(String[] args) {
        Date d1 = new Date(7, 1, 2025); // Istanzia it.oop.core.Date!

        // 2. Per istanziare la classe Date del package it.oop.ui,
        // è obbligatorio usare il Fully Qualified Name (FQN):
        it.oop.ui.Date d = new it.oop.ui.Date(12345);
        System.out.println(d.time);      // OK: 'time' è public
        // System.out.println(d.start);  // COMPILE-TIME ERROR: 'start' è private!
    }
}
\end{lstlisting}

\paragraph{\texorpdfstring{3. Anatomia del
\href{file:///home/wally/Documenti/GitHub/Programmazione-II/25-26/Teoria-e-Esercitazioni/Codice-20260902/Lezione10/MavenDate/pom.xml}{\texttt{pom.xml}}
di
\texttt{MavenDate}}{3. Anatomia del pom.xml di MavenDate}}\label{anatomia-del-pom.xml-di-mavendate}

\begin{lstlisting}[language=XML]
<project ...>
    <modelVersion>4.0.0</modelVersion>
    <groupId>it.oop</groupId>
    <artifactId>MavenDate</artifactId>
    <version>1.0-SNAPSHOT</version>

    <properties>
        <maven.compiler.source>17</maven.compiler.source>
        <maven.compiler.target>17</maven.compiler.target>
        <project.build.sourceEncoding>UTF-8</project.build.sourceEncoding>
    </properties>

    <build>
        <plugins>
            <!-- Plugin per creare il JAR eseguibile con dipendenze -->
            <plugin>
                <groupId>org.apache.maven.plugins</groupId>
                <artifactId>maven-assembly-plugin</artifactId>
                <version>3.7.1</version>
                <configuration>
                    <descriptorRefs>
                        <descriptorRef>jar-with-dependencies</descriptorRef>
                    </descriptorRefs>
                    <archive>
                        <manifest>
                            <addClasspath>true</addClasspath>
                            <!-- Entrypoint del JAR -->
                            <mainClass>it.oop.ui.MainDate</mainClass>
                        </manifest>
                    </archive>
                </configuration>
                <executions>
                    <execution>
                        <id>assemble-all</id>
                        <phase>package</phase>
                        <goals>
                            <goal>single</goal>
                        </goals>
                    </execution>
                </executions>
            </plugin>
        </plugins>
    </build>
</project>
\end{lstlisting}

Questo file è il prototipo esatto che hai esteso in tutti i moduli del
tuo workspace \passthrough{\lstinline!esercizi-wally-25-26!}: -
\passthrough{\lstinline!compile!}: compila da
\passthrough{\lstinline!src/main/java!} verso
\passthrough{\lstinline!target/classes!}. -
\passthrough{\lstinline!test!}: compila ed esegue i test da
\passthrough{\lstinline!src/test/java!} verso
\passthrough{\lstinline!target/test-classes!}. -
\passthrough{\lstinline!package!}: genera il JAR auto-eseguibile con
\passthrough{\lstinline!META-INF/MANIFEST.MF!} configurato con
\passthrough{\lstinline!Main-Class: it.oop.ui.MainDate!}.

\begin{center}\rule{0.5\linewidth}{0.5pt}\end{center}

\subsubsection{3. Ciclo: Corrispondenza con i Tuoi Moduli
Workspace}\label{ciclo-corrispondenza-con-i-tuoi-moduli-workspace}

Nel tuo repository
\href{file:///home/wally/Documenti/GitHub/Programmazione-II/25-26/esercizi-wally-25-26}{\passthrough{\lstinline!esercizi-wally-25-26!}}:
- \textbf{\passthrough{\lstinline!es08!}}:
\href{file:///home/wally/Documenti/GitHub/Programmazione-II/25-26/esercizi-wally-25-26/es08/src/main/java/es08/Memory.java}{\passthrough{\lstinline!Memory.java!}}
implementa l'allocazione Stack vs Heap e il blocco
\passthrough{\lstinline!static \{ ... \}!}. -
\textbf{\passthrough{\lstinline!es10!}}:
\href{file:///home/wally/Documenti/GitHub/Programmazione-II/25-26/esercizi-wally-25-26/es10/src/main/java/es10/ImmutablePerson.java}{\passthrough{\lstinline!ImmutablePerson.java!}}
applica formalmente tutte le regole di immutabilità di T10
(\passthrough{\lstinline!final class!}, campi
\passthrough{\lstinline!final!}, validazione e assenza di setter).

\begin{center}\rule{0.5\linewidth}{0.5pt}\end{center}

\begin{studynote}[Nota di Studio] Con la \textbf{Lezione 10} abbiamo completato l'intero
quadro su memoria, call stack, tipi wrapper e la transizione a Maven.

Il prossimo blocco è la \textbf{Lezione 11 / Slide T11: \emph{Packages
and Class Visibility}}: - Tassonomia gerarchica dei package e
convenzioni DNS invertite (\passthrough{\lstinline!it.univr...!}). -
Regole di visibilità e accessibilità inter-package: membri
\passthrough{\lstinline!public!}, \passthrough{\lstinline!protected!},
package-private e \passthrough{\lstinline!private!}. - Evoluzione di
\passthrough{\lstinline!MavenDate!}: - Refactoring dell'algoritmo
Gregoriano: introduzione definitiva del \textbf{calcolo dell'anno
bisestile (\passthrough{\lstinline!isLeapYear!})} in
\passthrough{\lstinline!Date.java!}! - Integrazione di
\passthrough{\lstinline!isLeapYear(year)!} in
\passthrough{\lstinline!daysPerMonth(month, year)!} e fine dei mesi
errati per Febbraio!
\end{studynote}

\textbf{Dammi conferma per procedere alla Lezione 11 / T11!}


\end{document}
